\documentclass[11pt,a4paper]{article}
\usepackage[utf8]{inputenc}
\usepackage[vietnamese]{babel}
\usepackage{amsmath,amssymb}
\usepackage{geometry}
\geometry{a4paper, total={180mm,260mm}, left=15mm, right=15mm, top=18mm, bottom=18mm}

\usepackage{tikz}
\usetikzlibrary{shapes.geometric, arrows.meta, positioning, calc}
\tikzset{
  arr/.style={-{Stealth[scale=0.8]}},
  every node/.append style={align=center}
}
\usepackage{tcolorbox}
\tcbuselibrary{skins,breakable}
\usepackage{fontawesome}
\usepackage{fancyhdr}
\usepackage{enumitem}

\pagestyle{fancy}
\fancyhf{}
\renewcommand{\headrulewidth}{0.5pt}
\renewcommand{\footrulewidth}{0.5pt}
\fancyhead[L]{\small\bfseries\color{violet!85!black}\faBook\ SỔ TAY KIẾN THỨC TOÁN 2}
\fancyhead[R]{\small\bfseries\color{violet!85!black}Lớp Toán Thầy Lĩnh -- loptoanthaylinh.vn}
\fancyfoot[C]{\small\color{gray}-- Sổ tay kiến thức: Bài 41. Phép chia --}

\setlength{\parindent}{0pt}
\setlength{\parskip}{0pt}

\begin{document}

% --- BANNER BÀI HỌC ---
\begin{tcolorbox}[
  enhanced,
  colback=violet!8!white,
  colframe=violet!70!black,
  arc=3mm, boxrule=1.2pt,
  left=3mm, right=3mm, top=2mm, bottom=2mm,
  halign=center
]
  {\LARGE\bfseries\color{violet!90!black}\faBookmark\ BÀI 41: PHÉP CHIA}\\[1mm]
  {\itshape\footnotesize\color{black!75} Hình thành phép chia từ tình huống chia đều và quan hệ đảo ngược với phép nhân trên 1 trang A4}
\end{tcolorbox}

\vspace{-2mm}

\noindent
\begin{minipage}[t]{0.49\textwidth}\vspace{0pt}
\begin{tcolorbox}[
  enhanced,
  colback=violet!4!white,
  colframe=violet!75!black,
  arc=2mm, boxrule=0.8pt,
  left=2mm, right=2mm, top=1.8mm, bottom=1.8mm,
  title={\footnotesize\bfseries\color{white}\faKey\ 1. Khái niệm phép chia}
]
{\footnotesize
\begin{itemize}\setlength{\itemsep}{1pt}\setlength{\leftskip}{-2.5mm}
  \item \textbf{Phép chia} xuất phát từ tình huống \textbf{chia đều} một số lượng thành các phần bằng nhau.
  \item \textbf{Dạng 1: Chia thành các phần bằng nhau}\\
  Có $6$ cái kẹo chia đều cho $3$ bạn. Mỗi bạn được:
  $$6 : 3 = \mathbf{2}\text{ (cái kẹo)}$$
  \item \textbf{Dạng 2: Chia theo nhóm}\\
  Có $6$ cái kẹo, chia cho mỗi bạn $2$ cái. Số bạn được kẹo là:
  $$6 : 2 = \mathbf{3}\text{ (bạn)}$$
  \item Dấu $:$ đọc là \textbf{"chia"}.
\end{itemize}
}
\end{tcolorbox}
\end{minipage}
\hfill
\begin{minipage}[t]{0.49\textwidth}\vspace{0pt}
\begin{tcolorbox}[
  enhanced,
  colback=teal!4!white,
  colframe=teal!75!black,
  arc=2mm, boxrule=0.8pt,
  left=2mm, right=2mm, top=1.8mm, bottom=1.8mm,
  title={\footnotesize\bfseries\color{white}\faSliders\ 2. Quan hệ giữa Phép nhân và Phép chia}
]
{\footnotesize
\begin{itemize}\setlength{\itemsep}{1pt}\setlength{\leftskip}{-2.5mm}
  \item Từ một phép tính nhân, ta luôn lập được \textbf{hai phép tính chia tương ứng}:
  $$\mathbf{a \times b = c \implies c : a = b \quad\text{và}\quad c : b = a}$$
  \item \textbf{Ví dụ:}
  Từ $2 \times 4 = 8$, ta có:
  $$8 : 2 = 4 \quad\text{và}\quad 8 : 4 = 2$$
  \item Phép chia là phép toán \textbf{ngược} của phép nhân.
\end{itemize}
}
\end{tcolorbox}
\end{minipage}

\vspace{-1.5mm}
\begin{tcolorbox}[
  enhanced,
  colback=orange!4!white,
  colframe=orange!80!black,
  arc=2mm, boxrule=0.8pt,
  left=2mm, right=2mm, top=1.5mm, bottom=1.5mm,
  title={\footnotesize\bfseries\color{white}\faSitemap\ 3. Sơ đồ tư duy liên kết: Một phép nhân tạo hai phép chia}
]
\centering
\begin{tikzpicture}[x=1.5cm, y=0.8cm]
  \node[rectangle, rounded corners=2pt, draw=violet!80!black, fill=violet!20, font=\scriptsize\bfseries] (mul) at (0,0) {
    PHÉP NHÂN\\$3 \times 5 = 15$
  };
  \node[rectangle, rounded corners=2pt, draw=teal!80!black, fill=teal!15, font=\scriptsize\bfseries, above right=0.1cm and 1.5cm of mul] (d1) {
    PHÉP CHIA 1\\$15 : 3 = 5$
  };
  \node[rectangle, rounded corners=2pt, draw=orange!80!black, fill=orange!15, font=\scriptsize\bfseries, below right=0.1cm and 1.5cm of mul] (d2) {
    PHÉP CHIA 2\\$15 : 5 = 3$
  };

  \draw[arr, thick, teal!80!black, ->] (mul) -- (d1);
  \draw[arr, thick, orange!80!black, ->] (mul) -- (d2);
\end{tikzpicture}
\end{tcolorbox}

\vspace{-1.5mm}
\noindent
\begin{minipage}[t]{0.49\textwidth}\vspace{0pt}
\begin{tcolorbox}[
  enhanced,
  colback=green!4!white,
  colframe=green!60!black,
  arc=2mm, boxrule=0.8pt,
  left=2mm, right=2mm, top=1.8mm, bottom=1.8mm,
  title={\footnotesize\bfseries\color{white}\faCheckCircle\ 4. Ví dụ thực hành chuẩn mực}
]
{\footnotesize
\textbf{Từ phép nhân viết phép chia:}\\
Cho phép nhân: $2 \times 6 = 12$. Viết hai phép chia tương ứng:\\
$\implies 12 : 2 = \mathbf{6}$.\\
$\implies 12 : 6 = \mathbf{2}$.\\
Cho phép nhân: $5 \times 4 = 20$:\\
$\implies 20 : 5 = \mathbf{4}$ và $20 : 4 = \mathbf{5}$.
}
\end{tcolorbox}
\end{minipage}
\hfill
\begin{minipage}[t]{0.49\textwidth}\vspace{0pt}
\begin{tcolorbox}[
  enhanced,
  colback=red!4!white,
  colframe=red!75!black,
  arc=2mm, boxrule=0.8pt,
  left=2mm, right=2mm, top=1.8mm, bottom=1.8mm,
  title={\footnotesize\bfseries\color{white}\faExclamationTriangle\ 5. Mẹo tránh bẫy trắc nghiệm}
]
{\footnotesize
\begin{itemize}\setlength{\itemsep}{1pt}\setlength{\leftskip}{-2.5mm}
  \item \textbf{Bản chất là "chia đều":}\\
  Chỉ làm phép chia khi đề bài nêu rõ các phần được chia \textbf{bằng nhau}.
  \item Để tìm kết quả phép chia, ta luôn nhẩm lại bảng nhân tương ứng.
\end{itemize}
}
\end{tcolorbox}
\end{minipage}


\end{document}
